% Drawing 69/01-15 -- folder 69, batch 01, archive page 20 (PDF page 21).
% Transcription: à gauche du bloc biffé, deux segments issus d'un même point, aux extrémités marquées $\xi_1 \to \sigma\xi_1$ et $\xi_2$, $\sigma\xi_2$, joints par de petits traits
% Kind: two segments from one point. Method 2 (schematic redraw from the facsimile). Status: ok.
% Notes: From a common point, segments to xi_1 and to xi_2, a third stroke between them; short strokes join xi_1 to sigma xi_1 and xi_2 to sigma xi_2.
% Coordinates: 1 unit = 1 mm = 1/100 of the crop width (bbox in page fractions: [0.1, 0.48, 0.34, 0.58]).
% Redrawn by Claude (claude-opus-5-5), issue #73, 2026-11-07; not checked by a human; not compiled (no TeX in the session).
% Rights: reproduces part of the work -- behind RENDER_TRANSCRIPTS, never published by itself.
\documentclass[tikz,border=3pt]{standalone}
\usepackage{amsmath,amssymb}
\usetikzlibrary{patterns,arrows.meta}
\tikzset{pencil/.style={black!45}, struck/.style={black!35}, mark/.style={circle,draw,inner sep=1.3pt}, pt/.style={circle,fill,inner sep=0.9pt}, opt/.style={circle,draw,fill=white,inner sep=1.1pt}}
\begin{document}
\begin{tikzpicture}[x=1mm,y=1mm,line cap=round,line join=round,font=\scriptsize]
\draw (21,-46) -- (27,-15);
\draw (21,-46) -- (55,-25);
\draw (21,-46) -- (60,-31);
\draw (27,-15) -- (36,-16);
\draw (55,-25) -- (60,-31);
\node[inner sep=1pt] at (25,-10) {\scriptsize $\xi_1$};
\node[inner sep=1pt] at (41,-14) {\scriptsize $\sigma\xi_1$};
\node[inner sep=1pt] at (56,-21) {\scriptsize $\xi_2$};
\node[inner sep=1pt] at (67,-31) {\scriptsize $\sigma\xi_2$};
\end{tikzpicture}
\end{document}
